\documentclass[12pt]{article}
\usepackage[utf8]{inputenc}
\usepackage{amsmath, amssymb, amsthm}
\usepackage{geometry}

% Set margins to match a standard homework document
\geometry{a4paper, margin=1in}

\title{Challenge Problem Set 2}
\author{Math 292 Spring 2026 \\ Rutgers University}
\date{}

\begin{document}

\maketitle{}

\noindent{}\textbf{Exercise 1:}

1. We can differentiate our solution x(t) to get
\begin{align*}
    \mathbf{x}''(t) = \sum_{j=0}^4{x''_j(t)}
\end{align*}

2. Thus we have for $x_0$ that:
\[
    x''_0(t) = -Kx_0(t)
\]

3. And we have for $j \in\{1,2,3,4\}$ that
\[
    x''_j(t) = -Kx_j(t) + f_j(t)
\]

4. Combining the two previous equations we have
\begin{align*}
    \mathbf{x}''(t) & = -Kx_0(t) + \sum_{j=1}^4(-Kx_j(t) + f_j(t))    \\
                    & = -K\sum_{i=0}^4(x_i(t)) + \sum_{j=1}^4(f_j(t)) \\
                    & = -K\mathbf{x}(t) + \mathbf{f}(t)
\end{align*}

5. This is exactly the original differential equation.

6. The initial conditions also check out
\begin{align*}
    \mathbf{x}(0)  & = \sum_{j=0}^4(x_j(t))                   \\
                   & = (1,2 ) + (0,0) + (0,0) + (0,0) + (0,0) \\
                   & = (1,2)                                  \\\\
    \mathbf{x}'(0) & = \sum_{j=0}^4 (x'_j(t))                 \\
                   & = (1,1) + (0,0) + (0,0) + (0,0) + (0,0)  \\
                   & = (1,1)
\end{align*}

\newpage{}

\noindent{}\textbf{Exercise 2:}

1. We begin by diagonalizing K in order the uncouple the system. To do this, we first need to solve for the eigenvalues using the characteristic polynomial of K.
\begin{align*}
    det\begin{pmatrix} 3 - \lambda{} & 2 \\ 2 & 6 - \lambda{}\end{pmatrix} & = (3-\lambda{})(6-\lambda{}) - 4 \\
                                                                           & = \lambda^{2} -9\lambda{} + 14   \\
                                                                           & = (\lambda{} - 7)(\lambda{} - 2) \\
                                                                           & = 0
\end{align*}

2. We now have that the eigenvalues are
\begin{align*}
    \lambda_{1} = 7 , ~ \lambda_{2} = 2
\end{align*}

3. With this, we can solve for the eigenvectors of K
\begin{itemize}
    \item $\lambda_{1} = 7$:
          \begin{align*}
              \begin{bmatrix} -4 & 2 \\ 2 & -1 \end{bmatrix}\mathbf{v}_1 & = \mathbf{0}                                              \\
              \mathbf{v}_1                                               & = \frac{1}{\sqrt{5}}\begin{bmatrix} 1 \\  2 \end{bmatrix}
          \end{align*}
    \item $\lambda_{1} = 2$:
          \begin{align*}
              \begin{bmatrix} 1 & 2 \\ 2 & 4 \end{bmatrix}\mathbf{v}_2 & = \mathbf{0}                                               \\
              \mathbf{v}_2                                             & = \frac{1}{\sqrt{5}}\begin{bmatrix} -2 \\  1 \end{bmatrix}
          \end{align*}
\end{itemize}

4. Thus, using the diagonalization of K, we have
\begin{align*}
    \mathbf{x}''(t) & = -VSV^{-1} \mathbf{x}(t)                                                                                                                                                \\
                    & = -\frac{1}{5} \begin{bmatrix} 1 & 2 \\ 2 & - 1 \end{bmatrix}  \begin{bmatrix} 7 & 0 \\ 0 & 2 \end{bmatrix}  \begin{bmatrix} 1 & 2 \\ -2 & 1 \end{bmatrix} \mathbf{x}(t)
\end{align*}

5. We then define a new variable y
\begin{align*}
    \mathbf{y} & = V^{-1}\mathbf{x}(t) = \frac{1}{\sqrt{5}} \begin{bmatrix} 1x_1+2x_2 \\ -2x_1+1x_2\end{bmatrix} \\
\end{align*}

6. Rearranging the equation we have
\begin{align*}
    \mathbf{x}''(t)       & = -VS\mathbf{y(t)} \\
    V^{-1}\mathbf{x}''(t) & = -S\mathbf{y(t)}  \\
    \mathbf{y}''          & = -S\mathbf{y}
\end{align*}

7. We have the following system
\begin{align*}
    y''_1 & = -7y_1 \\
    y''_2 & = -2y_2
\end{align*}

8. This is the same as the solutions of the previous workshop, thus we have
\begin{align*}
    y_1(t)  & = c_1\cos(\sqrt{7}t) + c_2\sin(\sqrt{7}t)                \\
    y_2(t)  & = c_3\cos(\sqrt{2}t) + c_4\sin(\sqrt{2}t)                \\
    y_1'(t) & = \sqrt{7}c_2\cos(\sqrt{7}t) -\sqrt{7}c_1\sin(\sqrt{7}t) \\
    y_2'(t) & = \sqrt{2}c_4\cos(\sqrt{2}t) -\sqrt{2}c_3\sin(\sqrt{2}t) \\
            & \text{with } c_1, c_2, c_3, c_4 \in{\mathbb{R}}
\end{align*}

9. At $t=0$ this would turn into
\begin{align*}
    y_1(0)  & = c_1         \\
    y_2(0)  & = c_3         \\
    y'_1(0) & = \sqrt{7}c_2 \\
    y'_2(0) & = \sqrt{2}c_4 \\
\end{align*}

9. We use the change of basis equation and initial conditions to solve for our constants
\begin{align*}
    \mathbf{y}(0) & = V^{-1}\mathbf{x}(0)                                                                                  \\
                  & = \frac{1}{\sqrt{5}}\begin{pmatrix} 1 & 2 \\ -2 & 1 \end{pmatrix} \begin{pmatrix} 1 \\ 2 \end{pmatrix} \\
                  & = \frac{1}{\sqrt{5}}\begin{pmatrix} 5 \\ 0  \end{pmatrix}
\end{align*}

10. We have our first two constants
\begin{align*}
    c_1 & = \sqrt{5} \\
    c_3 & = 0
\end{align*}

11. Now, we repeat with y'
\begin{align*}
    \mathbf{y}'(0) & = V^{-1}\mathbf{x}'(0)                                                                                 \\
                   & = \frac{1}{\sqrt{5}}\begin{pmatrix} 1 & 2 \\ -2 & 1 \end{pmatrix} \begin{pmatrix} 1 \\ 1 \end{pmatrix} \\
                   & = \frac{1}{\sqrt{5}}\begin{pmatrix} 3 \\ -1  \end{pmatrix}
\end{align*}

12. Solving the system we have
\begin{align*}
    c_2 & = \frac{3}{\sqrt{35}}  \\
    c_4 & = -\frac{1}{\sqrt{10}}
\end{align*}

13. Finally, we have the final solution:
\begin{align*}
    \mathbf{x}_0(t) & = V\mathbf{y}(t)                                                                                                                                                             \\
                    & = y_1(t)\mathbf{v_1} + y_2(t)\mathbf{v_2}                                                                                                                                    \\
                    & = (\cos(\sqrt{7}t) + \frac{3}{5\sqrt{7}}{\sin(\sqrt{7}t)}) \begin{bmatrix} 1 \\ 2 \end{bmatrix} - (\frac{1}{5\sqrt{2}}\sin(\sqrt{2}t)) \begin{bmatrix} -2 \\ 1 \end{bmatrix}
\end{align*}

14. Separating $\mathbf{x}_0$ into it's components we have
\begin{align*}
    x_1(t) & = \cos(\sqrt{7}t) + \frac{3}{5\sqrt{7}}sin(\sqrt{7}t) + \frac{2}{5\sqrt{2}}\sin(\sqrt{2}t)  \\
    x_2(t) & = 2\cos(\sqrt{7}t) + \frac{6}{5\sqrt{7}}sin(\sqrt{7}t) - \frac{1}{5\sqrt{2}}\sin(\sqrt{2}t)
\end{align*}

\noindent{}\textbf{Exercise 3:}

Now we have that $\mathbf{x}_1''(t) = -K\mathbf{x}_1(t) + \mathbf{f}_1(t)$,
but we can again use the diagonalization of K from the previous part to solve for the solution,
\begin{align*}
    \mathbf{x}_1''(t) & = -K\mathbf{x}_1(t) + \mathbf{f}_1(t)        \\
                      & = -VSV^{-1}\mathbf{x}_1(t) = \mathbf{f}_1(t)
\end{align*}

We let $\mathbf{y}(t) = V^{-1}\mathbf{x}(t)$ and multiply both sides by $V^{-1}$
\begin{align*}
    \mathbf{y}''(t) + S\mathbf{y}(t) = V^{-1}\mathbf{f}_1(t)
\end{align*}

We are left with the following system
\begin{align*}
    y_1''(t) + 7y_1 & = \frac{7}{\sqrt{5}} \cos(\omega_{1}t) \\
    y_2''(t) + 2y_2 & = \frac{1}{\sqrt{5}} \cos(\omega_{1}t)
\end{align*}

Similar to the last part, we guess the solution to y
\begin{align*}
    y_1(t) & = A_1\cos(\sqrt{7}t) + B_1\sin(\sqrt{7}t) + C_1cos(\omega_{1}t) \\
    y_2(t) & = A_2\cos(\sqrt{2}t) + B_2\sin(\sqrt{2}t) + C_2cos(\omega_{1}t)
\end{align*}

Then, taking the first and second derivative we have
\begin{align*}
    y'_1(t)  & = -\sqrt{7}A_1\sin(\sqrt{7}t) + \sqrt{7}B_1\cos(\sqrt{7}t) - \omega_{1} C_1 sin(\omega_{1}t) \\
    y'_2(t)  & = -\sqrt{2}A_2\sin(\sqrt{2}t) + \sqrt{2}B_2\cos(\sqrt{2}t) - \omega_{1} C_2 sin(\omega_{1}t) \\\\
    y''_1(t) & = -7A_1\cos(\sqrt{7}t) - 7B_1\sin(\sqrt{7}t) - \omega_{1}^2 C_1cos(\omega_{1}t)              \\
    y''_2(t) & = -2A_2\cos(\sqrt{2}t) - 2B_2\sin(\sqrt{2}t) - \omega_{1}^2 C_2cos(\omega_{1}t)
\end{align*}

We solve for $C_1$ and $C_2$ by plugging into the first system we solved for
\begin{itemize}
    \item For $C_1$:
          \begin{align*}
              -\omega_1^2 C_1 \cos(\omega_1 t) + 7C_1cos(\omega_1t) & = \frac{7}{\sqrt{5}}cos(\omega_1t) \\
              C_1(7-\omega_1^2)                                     & = \frac{7}{\sqrt{5}}               \\
              C_1                                                   & = \frac{7}{\sqrt{5}(7-\omega_1^2)}
          \end{align*}
    \item For $C_2$:
          \begin{align*}
              -\omega_1^2 C_2 \cos(\omega_1 t) + 2C_2cos(\omega_1t) & = \frac{1}{\sqrt{5}}cos(\omega_1t) \\
              C_2(2-\omega_1^2)                                     & = \frac{1}{\sqrt{5}}               \\
              C_2                                                   & = \frac{1}{\sqrt{5}(2-\omega_1^2)}
          \end{align*}
\end{itemize}

Next, we solve for $A_1$ and $A_2$ with the initial condition $\mathbf{x}(0) = 0 \implies{} \mathbf{y}(0) = 0$
\begin{itemize}
    \item For $A_1$:
          \begin{align*}
              y_1(0) & = A_1\cos(0) + B_1\sin(0) + C_1cos(0) \\
              A_1    & = -C_1                                \\
                     & = -\frac{7}{\sqrt{5}(7-\omega_1^2)}
          \end{align*}
    \item For $A_2$:
          \begin{align*}
              y_2(0) & = A_2\cos(0) + B_2\sin(0) + C_2cos(0) \\
              A_2    & = -C_2                                \\
                     & = -\frac{1}{\sqrt{5}(2-\omega_1^2)}
          \end{align*}
\end{itemize}

We repeat this again for $B_1$ and $B_2$ using the initial condition $\mathbf{x}'(0) = 0 \implies{} \mathbf{y}'(0) = 0$
\begin{itemize}
    \item For $B_1$:
          \begin{align*}
              y'_1(0) & = -\sqrt{7}A_1\sin(0) + \sqrt{7}B_1\cos(0) - \omega_{1} C_1 sin(0) \\
              B_1     & = 0
          \end{align*}
    \item For $B_2$:
          \begin{align*}
              y'_2(0) & = -\sqrt{2}A_2\sin(0) + \sqrt{2}B_2\cos(0) - \omega_{1} C_2 sin(0) \\
              B_2     & = 0
          \end{align*}
\end{itemize}

Notice that $B_1$ and $B_2$ will always be 0 for $x_i ~|~ i = 1,2,3,4$ since only the middle term will not go to zero when t=0.


Thus, we have our solution for $\mathbf{y}(t)$, where $\omega_{i}(t) \neq{\sqrt{7}},\sqrt{2} ~|~ i = 1,2,3,4$
\begin{align*}
    \mathbf{y}(t) = \begin{bmatrix}
                        \dfrac{7}{\sqrt{5}(7-\omega_1^2)}(\cos(\omega_{1}t) -\cos(\sqrt{7}t)) \\
                        \dfrac{1}{\sqrt{5}(2-\omega_1^2)}(\cos(\omega_{1}t) -\cos(\sqrt{2}t))
                    \end{bmatrix}
\end{align*}

Converting back to our original basis with $\mathbf{x}(t) = V\mathbf{y}(t)$, we get
\begin{align*}
    \mathbf{x_1}(t) & = \begin{bmatrix}
                            \dfrac{7}{5(7-\omega_1^2)}(\cos(\omega_{1}t) -\cos(\sqrt{7}t)) - \dfrac{2}{5(2-\omega_1^2)}(\cos(\omega_{1}t) -\cos(\sqrt{2}t)) \\
                            \dfrac{14}{5(7-\omega_1^2)}(\cos(\omega_{1}t) -\cos(\sqrt{7}t)) + \dfrac{1}{5(2-\omega_1^2)}(\cos(\omega_{1}t) -\cos(\sqrt{2}t))
                        \end{bmatrix}
\end{align*}

However, we must also consider the case where $\omega_1 = \pm{}\sqrt{7}$ or $\omega_1 = \pm{}\sqrt{2}$ or when the system is in "resonance". Choosing the positive values, we have

\begin{itemize}
    \item $\omega_1 = {\sqrt{7}}$:

          We take the limit of $y_1(t)$ as the $\omega_1$ approaches the natural frequency, via L'Hopital's Rule
          \begin{align*}
              \lim_{\omega_1 \rightarrow{} \sqrt{7}}  \dfrac{7}{\sqrt{5}(7-\omega_1^2)}(\cos(\omega_{1}t) -\cos(\sqrt{7}t)) & \Rightarrow{} \frac{-7t \sin(\omega_1t)}{-2\sqrt{5}\omega_1} \\
                                                                                                                            & \Rightarrow{}  \frac{\sqrt{7}}{2\sqrt{5}} t\sin{(\sqrt{7}t)}
          \end{align*}

          Then, we do the same for $y_2(t)$ by just substituting in the limit
          \begin{align*}
              \lim_{\omega_2 \rightarrow{} \sqrt{7}} \dfrac{1}{\sqrt{5}(2-\omega_1^2)}(\cos(\omega_{1}t) -\cos(\sqrt{2}t)) \Rightarrow{} \frac{1}{5\sqrt{5}} (\cos(\sqrt{2}t) -\cos(\sqrt{7}t))
          \end{align*}

          Thus, for the final solution, we have that
          \begin{align*}
              \mathbf{x_1}(t) & = \begin{bmatrix}
                                      \frac{\sqrt{7}}{10} t\sin{(\sqrt{7}t)} - \frac{2}{25} (\cos(\sqrt{2}t) -\cos(\sqrt{7}t)) \\
                                      \frac{\sqrt{7}}{5} t\sin{(\sqrt{7}t)} + \frac{1}{25} (\cos(\sqrt{2}t) -\cos(\sqrt{7}t))  \\
                                  \end{bmatrix}
          \end{align*}

    \item $\omega_1 = \sqrt{2}$

          We take the limit of $y_1(t)$ as the $\omega_1$ approaches the natural frequency, substituting for $\sqrt{2}$
          \begin{align*}
              \lim_{\omega_1 \rightarrow{} \sqrt{2}}  \dfrac{7}{\sqrt{5}(7-\omega_1^2)}(\cos(\omega_{1}t) -\cos(\sqrt{7}t)) & \Rightarrow{} \dfrac{7}{5\sqrt{5}}(\cos(\sqrt{2}t) -\cos(\sqrt{7}t))
          \end{align*}

          Then, we do the same for $y_2(t)$ via L'Hopital's Rule
          \begin{align*}
              \lim_{\omega_2 \rightarrow{} \sqrt{2}} \dfrac{\cos(\omega_{1}t) -\cos(\sqrt{2}t)}{\sqrt{5}(2-\omega_1^2)} & \Rightarrow{} \dfrac{-t \sin(\omega_1t)}{-2\sqrt{5}\omega_1} \\
                                                                                                                        & \Rightarrow{} \dfrac{1}{2\sqrt{10}} t\sin(\sqrt{2}t)
          \end{align*}

          Thus, for the final solution, we have that
          \begin{align*}
              \mathbf{x_1}(t) & = \begin{bmatrix}
                                      \frac{7}{25}(\cos(\sqrt{2}t) -\cos(\sqrt{7}t)) - \frac{1}{5\sqrt{2}} t\sin(\sqrt{2}t) \\
                                      \frac{14}{25}(\cos(\sqrt{2}t) -\cos(\sqrt{7}t)) + \frac{1}{10\sqrt{2}} t\sin(\sqrt{2}t)
                                  \end{bmatrix}
          \end{align*}
\end{itemize}

\noindent{}\textbf{Exercise 4:}

We follow a similar process, to the last part, the only different part is f(t). Thus, we start with the change of basis $\mathbf{y}(t) = V^{-1}\mathbf{x}(t)$
\begin{align*}
    \mathbf{y}''(t) + S\mathbf{y}(t) = V^{-1}\mathbf{f}_2(t)
\end{align*}

This gives us the following system
\begin{align*}
    y''_1(t) + 7y_1 & = -\frac{1}{\sqrt{5}} \sin(\omega_2t) \\
    y''_2(t) + 2y_2 & = -\frac{3}{\sqrt{5}} \sin(\omega_2t)
\end{align*}

We guess the solution of y to be
\begin{align*}
    y_1(t) & = A_1\cos(\sqrt{7}t) + B_1\sin(\sqrt{7}t) + C_1 \sin(\omega_{2}t) \\
    y_2(t) & = A_2\cos(\sqrt{2}t) + B_2\sin(\sqrt{2}t) + C_2 \sin(\omega_{2}t)
\end{align*}

Then, taking the first and second derivative we have
\begin{align*}
    y'_1(t)  & = -\sqrt{7}A_1\sin(\sqrt{7}t) + \sqrt{7}B_1\cos(\sqrt{7}t) + \omega_{2} C_1 \cos(\omega_{2}t) \\
    y'_2(t)  & = -\sqrt{2}A_2\sin(\sqrt{2}t) + \sqrt{2}B_2\cos(\sqrt{2}t) + \omega_{2} C_2 \cos(\omega_{2}t) \\\\
    y''_1(t) & = -7A_1\cos(\sqrt{7}t) - 7B_1\sin(\sqrt{7}t) - \omega_{2}^2 C_1 \sin(\omega_{2}t)             \\
    y''_2(t) & = -2A_2\cos(\sqrt{2}t) - 2B_2\sin(\sqrt{2}t) - \omega_{2}^2 C_2 \sin(\omega_{2}t)
\end{align*}

We solve for $C_1$ and $C_2$ by plugging into the first system we solved for
\begin{itemize}
    \item For $C_1$:
          \begin{align*}
              -\omega_2^2 C_1 \sin(\omega_2 t) + 7C_1 \sin(\omega_2t) & = -\frac{1}{\sqrt{5}}sin(\omega_2t) \\
              C_1(7-\omega_2^2)                                       & = -\frac{1}{\sqrt{5}}               \\
              C_1                                                     & = \frac{-1}{\sqrt{5}(7-\omega_2^2)}
          \end{align*}
    \item For $C_2$:
          \begin{align*}
              -\omega_2^2 C_2 \sin(\omega_2 t) + 2C_2 \sin(\omega_2t) & = -\frac{3}{\sqrt{5}} \sin(\omega_2t) \\
              C_2(2-\omega_2^2)                                       & = -\frac{3}{\sqrt{5}}                 \\
              C_2                                                     & = \frac{-3}{\sqrt{5}(2-\omega_2^2)}
          \end{align*}
\end{itemize}

Next, we solve for $A_1$ and $A_2$ with the initial condition $\mathbf{x}(0) = 0 \implies{} \mathbf{y}(0) = 0$
\begin{itemize}
    \item For $A_1$:
          \begin{align*}
              y_1(0) & = A_1\cos(\sqrt{7}t) + B_1\sin(0) + C_1 \sin(0) \\
              A_1    & = 0
          \end{align*}
    \item For $A_2$:
          \begin{align*}
              y_2(0) & = A_2\cos(0) + B_2\sin(0) + C_2 \sin(0) \\
              A_2    & = 0
          \end{align*}
\end{itemize}

From this and the previous exercise, we can see that when $\mathbf{f}(t)$ is a multiple of $\cos(t)$, then the solution of y reduces to
\begin{align*}
    \mathbf{y} = \begin{bmatrix}
                     A_1\cos(\sqrt{7}t) + C_1cos(\omega_{i}t) \\
                     A_2\cos(\sqrt{2}t) + C_2cos(\omega_{i}t)
                 \end{bmatrix} ~|~ i = 1,4
\end{align*}
and when $\mathbf{f}(t)$ is a multiple of $\sin(t)$, then the solution of y reduces to
\begin{align*}
    \mathbf{y} = \begin{bmatrix}
                     B_1\sin(\sqrt{7}t) + C_1 \sin(\omega_{i}t) \\
                     B_2\sin(\sqrt{2}t) + C_2 \sin(\omega_{i}t)
                 \end{bmatrix} ~|~ i = 2,3
\end{align*}

We can reuse this fact in the next two exercises.

Now, we solve for $B_1$ and $B_2$ using the initial condition $\mathbf{x}'(0) = 0 \implies{} \mathbf{y}'(0) = 0$
\begin{itemize}
    \item For $B_1$:
          \begin{align*}
              y'_1(0) & = -\sqrt{7}A_1\sin(0) + \sqrt{7}B_1\cos(0) + \omega_{2} C_1 \cos(0) \\
              B_1     & = -\frac{\omega_{2} C_1}{\sqrt{7}}                                  \\
                      & = \frac{\omega_{2}}{\sqrt{35}(7-\omega_2^2)}
          \end{align*}
    \item For $B_2$:
          \begin{align*}
              y'_2(0) & = -\sqrt{2}A_2\sin(0) + \sqrt{2}B_2\cos(0) + \omega_{2} C_2 \cos(0) \\
              B_2     & = -\frac{\omega_{2} C_2}{\sqrt{2}}                                  \\
                      & = \frac{3\omega_{2}}{\sqrt{10}(2-\omega_2^2)}
          \end{align*}
\end{itemize}

Thus, we have our solution for $\mathbf{y}(t)$,
\begin{align*}
    \mathbf{y} = \begin{bmatrix}
                     \frac{1}{\sqrt{5}(7-\omega_2^2)}(\frac{\omega_2}{\sqrt{7}}\sin(\sqrt{7}t) - \sin(\omega_{2}t)) \\
                     \frac{3}{\sqrt{5}(2-\omega_2^2)}(\frac{\omega_2}{\sqrt{2}}\sin(\sqrt{2}t) - \sin(\omega_{2}t))
                 \end{bmatrix}
\end{align*}

Converting back to our original basis with $\mathbf{x}(t) = V\mathbf{y}(t)$, we get
\begin{align*}
    \mathbf{x_2}(t) & = \begin{bmatrix}
                            \frac{1}{5(7-\omega_2^2)}(\frac{\omega_2}{\sqrt{7}}\sin(\sqrt{7}t) - \sin(\omega_{2}t)) - \frac{6}{5(2-\omega_2^2)}(\frac{\omega_2}{\sqrt{2}}\sin(\sqrt{2}t) - \sin(\omega_{2}t)) \\
                            \frac{2}{5(7-\omega_2^2)}(\frac{\omega_2}{\sqrt{7}}\sin(\sqrt{7}t) - \sin(\omega_{2}t)) + \frac{3}{5(2-\omega_2^2)}(\frac{\omega_2}{\sqrt{2}}\sin(\sqrt{2}t) - \sin(\omega_{2}t))
                        \end{bmatrix}
\end{align*}

\noindent{}\textbf{Exercise 5:}

We follow a similar process, to the last part, the only different part is f(t). Thus, we start with the change of basis $\mathbf{y}(t) = V^{-1}\mathbf{x}(t)$
\begin{align*}
    \mathbf{y}''(t) + S\mathbf{y}(t) = V^{-1}\mathbf{f}_3(t)
\end{align*}

This gives us the following system
\begin{align*}
    y''_1(t) + 7y_1 & = \frac{1}{\sqrt{5}} \sin(\omega_3t)  \\
    y''_2(t) + 2y_2 & = -\frac{7}{\sqrt{5}} \sin(\omega_3t)
\end{align*}

We guess the solution of y to be
\begin{align*}
    y_1(t) & =  B_1 \sin(\sqrt{7}t) + C_1 \sin(\omega_{3}t) \\
    y_2(t) & =  B_2 \sin(\sqrt{2}t) + C_2 \sin(\omega_{3}t)
\end{align*}

Then, taking the first and second derivative we have
\begin{align*}
    y'_1(t)  & =  \sqrt{7}B_1\cos(\sqrt{7}t) + \omega_{3} C_1 \cos(\omega_{3}t) \\
    y'_2(t)  & =  \sqrt{2}B_2\cos(\sqrt{2}t) + \omega_{3} C_2 \cos(\omega_{3}t) \\\\
    y''_1(t) & =  - 7B_1 \sin(\sqrt{7}t) - \omega_{3}^2 C_1 \sin(\omega_{3}t)   \\
    y''_2(t) & =  - 2B_2 \sin(\sqrt{2}t) - \omega_{3}^2 C_2 \sin(\omega_{3}t)
\end{align*}

We solve for $C_1$ and $C_2$ by plugging into the first system we solved for
\begin{itemize}
    \item For $C_1$:
          \begin{align*}
              -\omega_3^2 C_1 \sin(\omega_3 t) + 7C_1 \sin(\omega_3t) & = \frac{1}{\sqrt{5}} \sin(\omega_3t) \\
              C_1(7-\omega_3^2)                                       & = \frac{1}{\sqrt{5}}                 \\
              C_1                                                     & = \frac{1}{\sqrt{5}(7-\omega_3^2)}
          \end{align*}
    \item For $C_2$:
          \begin{align*}
              -\omega_3^2 C_2 \sin(\omega_3 t) + 2C_2 \sin(\omega_3t) & = -\frac{7}{\sqrt{5}} \sin(\omega_3t) \\
              C_2(2-\omega_3^2)                                       & = -\frac{7}{\sqrt{5}}                 \\
              C_2                                                     & = \frac{-7}{\sqrt{5}(2-\omega_3^2)}
          \end{align*}
\end{itemize}

Now, we solve for $B_1$ and $B_2$ using the initial condition $\mathbf{x}'(0) = 0 \implies{} \mathbf{y}'(0) = 0$
\begin{itemize}
    \item For $B_1$:
          \begin{align*}
              B_1 & = -\frac{\omega_{3} C_1}{\sqrt{7}}            \\
                  & = \frac{-\omega_{3}}{\sqrt{35}(7-\omega_3^2)}
          \end{align*}
    \item For $B_2$:
          \begin{align*}
              B_2 & = -\frac{\omega_{3} C_2}{\sqrt{2}}            \\
                  & = \frac{7\omega_{3}}{\sqrt{10}(2-\omega_3^2)}
          \end{align*}
\end{itemize}

Thus, we have our solution for $\mathbf{y}(t)$,
\begin{align*}
    \mathbf{y} = \begin{bmatrix}
                     \frac{1}{\sqrt{5}(7-\omega_3^2)}(\sin(\omega_{3}t) - \frac{\omega_3}{\sqrt{7}}\sin(\sqrt{7}t)) \\
                     \frac{-7}{\sqrt{5}(2-\omega_3^2)}(\sin(\omega_{3}t) - \frac{\omega_3}{\sqrt{2}}\sin(\sqrt{2}t))
                 \end{bmatrix}
\end{align*}

Converting back to our original basis with $\mathbf{x}(t) = V\mathbf{y}(t)$, we get
\begin{align*}
    \mathbf{x_3}(t) & = \begin{bmatrix}
                            \frac{1}{5(7-\omega_3^2)}(\sin(\omega_{3}t) - \frac{\omega_3}{\sqrt{7}}\sin(\sqrt{7}t)) + \frac{14}{5(2-\omega_3^2)}(\sin(\omega_{3}t) - \frac{\omega_3}{\sqrt{2}}\sin(\sqrt{2}t)) \\
                            \frac{2}{5(7-\omega_3^2)}(\sin(\omega_{3}t) - \frac{\omega_3}{\sqrt{7}}\sin(\sqrt{7}t)) + \frac{-7}{5(2-\omega_3^2)}(\sin(\omega_{3}t) - \frac{\omega_3}{\sqrt{2}}\sin(\sqrt{2}t))
                        \end{bmatrix}
\end{align*}

\noindent{}\textbf{Exercise 6:}

We follow a similar process, to the last part, the only different part is f(t). Thus, we start with the change of basis $\mathbf{y}(t) = V^{-1}\mathbf{x}(t)$
\begin{align*}
    \mathbf{y}''(t) + S\mathbf{y}(t) = V^{-1}\mathbf{f}_4(t)
\end{align*}

This gives us the following system
\begin{align*}
    y''_1(t) + 7y_1 & = \frac{1}{\sqrt{5}} \cos(\omega_4t)  \\
    y''_2(t) + 2y_2 & = -\frac{2}{\sqrt{5}} \cos(\omega_4t)
\end{align*}

We guess the solution of y to be
\begin{align*}
    y_1(t) & = A_1\cos(\sqrt{7}t) +  C_1cos(\omega_{4}t) \\
    y_2(t) & = A_2\cos(\sqrt{2}t) +  C_2cos(\omega_{4}t)
\end{align*}

Then, taking the first and second derivative we have
\begin{align*}
    y'_1(t)  & = -\sqrt{7}A_1\sin(\sqrt{7}t)  - \omega_{4} C_1 sin(\omega_{4}t) \\
    y'_2(t)  & = -\sqrt{2}A_2\sin(\sqrt{2}t)  - \omega_{4} C_2 sin(\omega_{4}t) \\\\
    y''_1(t) & = -7A_1\cos(\sqrt{7}t)  - \omega_{4}^2 C_1cos(\omega_{4}t)       \\
    y''_2(t) & = -2A_2\cos(\sqrt{2}t)  - \omega_{4}^2 C_2cos(\omega_{4}t)
\end{align*}

We solve for $C_1$ and $C_2$ by plugging into the first system we solved for
\begin{itemize}
    \item For $C_1$:
          \begin{align*}
              C_1(7-\omega_4^2) & = \frac{1}{\sqrt{5}}               \\
              C_1               & = \frac{1}{\sqrt{5}(7-\omega_4^2)}
          \end{align*}
    \item For $C_2$:
          \begin{align*}
              C_2(2-\omega_4^2) & = -\frac{2}{\sqrt{5}}               \\
              C_2               & = \frac{-2}{\sqrt{5}(2-\omega_4^2)}
          \end{align*}
\end{itemize}

Next, we solve for $A_1$ and $A_2$ with the initial condition $\mathbf{x}(0) = 0 \implies{} \mathbf{y}(0) = 0$
\begin{itemize}
    \item For $A_1$:
          \begin{align*}
              y_1(0) & = A_1\cos(0) + B_1\sin(0) + C_1cos(0) \\
              A_1    & = -C_1                                \\
                     & = -\frac{1}{\sqrt{5}(7-\omega_4^2)}
          \end{align*}
    \item For $A_2$:
          \begin{align*}
              y_2(0) & = A_2\cos(0) + B_2\sin(0) + C_2cos(0) \\
              A_2    & = -C_2                                \\
                     & = -\frac{-2}{\sqrt{5}(2-\omega_4^2)}
          \end{align*}
\end{itemize}

Thus, we have our solution for $\mathbf{y}(t)$,
\begin{align*}
    \mathbf{y} = \begin{bmatrix}
                     \frac{1}{\sqrt{5}(7-\omega_4^2)}(\cos(\omega_{4}t) - \cos(\sqrt{7}t)) \\
                     \frac{-2}{\sqrt{5}(2-\omega_4^2)}(\cos(\omega_{4}t) - \cos(\sqrt{2}t))
                 \end{bmatrix}
\end{align*}

Converting back to our original basis with $\mathbf{x}(t) = V\mathbf{y}(t)$, we get
\begin{align*}
    \mathbf{x_4}(t) & = \begin{bmatrix}
                            \frac{1}{5(7-\omega_4^2)}(\cos(\omega_{4}t) - \cos(\sqrt{7}t)) + \frac{4}{5(2-\omega_4^2)}(\cos(\omega_{4}t) - \cos(\sqrt{2}t)) \\
                            \frac{2}{5(7-\omega_4^2)}(\cos(\omega_{4}t) - \cos(\sqrt{7}t)) - \frac{2}{5(2-\omega_4^2)}(\cos(\omega_{4}t) - \cos(\sqrt{2}t))
                        \end{bmatrix}
\end{align*}

\newpage{}

\noindent{}\textbf{Exercise 7:}

Combining our solutions we have the final solution for (0.1) and (0.2)
\begin{align*}
    \mathbf{x}(t)             & = (\frac{1}{\sqrt{5}} y_1(t)) \begin{bmatrix} 1 \\ 2 \end{bmatrix} + (\frac{1}{\sqrt{5}} y_2(t)) \begin{bmatrix} -2 \\ 1 \end{bmatrix} \\\\
    \frac{1}{\sqrt{5}} y_1(t) & = \cos(\sqrt{7}t) + \frac{3}{5\sqrt{7}}{\sin(\sqrt{7}t)}                                                                               \\
                              & + \dfrac{7}{5(7-\omega_1^2)}(\cos(\omega_{1}t) -\cos(\sqrt{7}t))                                                                       \\
                              & + \frac{1}{5(7-\omega_2^2)}(\frac{\omega_2}{\sqrt{7}}\sin(\sqrt{7}t) - \sin(\omega_{2}t))                                              \\
                              & + \frac{1}{5(7-\omega_3^2)}(\sin(\omega_{3}t) - \frac{\omega_3}{\sqrt{7}}\sin(\sqrt{7}t))                                              \\
                              & + \frac{1}{5(7-\omega_4^2)}(\cos(\omega_{4}t) - \cos(\sqrt{7}t))                                                                       \\
    \frac{1}{\sqrt{5}} y_2(t) & = \frac{-1}{5\sqrt{2}}\sin(\sqrt{2}t)                                                                                                  \\
                              & + \dfrac{1}{5(2-\omega_1^2)}(\cos(\omega_{1}t) -\cos(\sqrt{2}t))                                                                       \\
                              & +  \frac{3}{5(2-\omega_2^2)}(\frac{\omega_2}{\sqrt{2}}\sin(\sqrt{2}t) - \sin(\omega_{2}t))                                             \\
                              & +  \frac{-7}{5(2-\omega_3^2)}(\sin(\omega_{3}t) - \frac{\omega_3}{\sqrt{2}}\sin(\sqrt{2}t))                                            \\
                              & + \frac{-2}{5(2-\omega_4^2)}(\cos(\omega_{4}t) - \cos(\sqrt{2}t))
\end{align*}

Notice that the system will exhibit resonance when $\omega_i = \sqrt{7}$ or $ \omega_i = \sqrt{2} ~|~ i = 1,2,3,4$

\noindent{}\textbf{Exercise 8:}

The closer the $\omega{}$ values are to resonance, the closer the denominator will be to zero which will lead to the function to blow up.

Approximating, we have that
\begin{align*}
    \sqrt{2} & \approx{1.414} \\
    \sqrt{7} & \approx{2.646}
\end{align*}

And the closest $\omega_i ~|~ i = 1,2,3,4$ is $\omega_3 = 2.65$.

Therefore, the solution will still be well approximated by just keeping the $\mathbf{x_3}(t)$ term or that
\begin{align*}
    \mathbf{x}(t) & \approx \begin{bmatrix}
                                \frac{1}{5(7-\omega_3^2)}(\sin(\omega_{3}t) - \frac{\omega_3}{\sqrt{7}}\sin(\sqrt{7}t)) + \frac{14}{5(2-\omega_3^2)}(\sin(\omega_{3}t) - \frac{\omega_3}{\sqrt{2}}\sin(\sqrt{2}t)) \\
                                \frac{2}{5(7-\omega_3^2)}(\sin(\omega_{3}t) - \frac{\omega_3}{\sqrt{7}}\sin(\sqrt{7}t)) + \frac{-7}{5(2-\omega_3^2)}(\sin(\omega_{3}t) - \frac{\omega_3}{\sqrt{2}}\sin(\sqrt{2}t))
                            \end{bmatrix} \\
                  & \text{with } \omega_3 = 2.65
\end{align*}

\end{document}